\chapter{Complex Analysis in Banach Spaces}

\section{Holomorphic Functions}

\begin{defn}Let $X$ be a complex Banach space, let $U \subset
  \complexes$ be open and let $f : U \to X$ be a function.   We say
  that $f$ is \emph{holomorphic} at $z \in U$ if the limit
  \begin{align*}
    f^\prime(z) &= \lim_{h \to 0} \frac{f(z+h) - f(z)}{h}
  \end{align*}
  exists (note $h$ is assumed to be a complex number).   If $f : \complexes \to X$ is holomorphic at every 
$z\in \complexes$ then we say that $f$ is \emph{entire}.
\end{defn}

The first class of examples of holomorphic functions is complex analytic functions.
\begin{prop}\label{banach:ComplexAnalyticImpliesHolomorphic}Let $X$ be a complex Banach space, let $U \subset
  \complexes$ be open and let $f : U \to X$ be an analytic function, then $f$ is holomorphic at every $z \in U$.
\end{prop}
\begin{proof}
  Let $z_0 \in U$ and pick a $\rho>0$ and a sequence $a_n$ in $X$ such that the power series $S(z) = \sum_{n=0}^\infty a_n z^n$ has radius of convergence greater than or equal to $\rho$ and 
  \begin{align*}
    f(z) = S(z - z_0) = \sum_{n=0}^\infty a_n (z - z_0)^n \text{ for $\abs{z - z_0} < \rho$}
  \end{align*}
  Then by Proposition \ref{banach:DifferentiationOfPowerSeries} we know that derived power series $S^\prime(z) = \sum_{n=1}^\infty n a_n z^{n-1}$ has radius of convergence greater than or equal to $\rho$ and,
  \begin{align*}
    a_1 &= S^\prime(0) = \lim_{h \to 0} \frac{S(h) - S(0)}{h} = \lim_{h \to 0} \frac{f(z_0 + h) - f(z_0)}{h} = f^\prime(z_0)
  \end{align*}
  which shows $f$ is holomorphic at $z_0$.   Since the $z_0$ was arbitrary the result follows.
\end{proof}

\begin{examp}Let $a_n$ be a sequence in a complex Banach space $X$ and
  suppose the series
  $S(z) = \sum_{n=0}^\infty a_n z^n$ has radius of convergence $\rho>0$.
  Then by Proposition \ref{banach:ConvergentPowerSeriesAreAnalytic} we
  know that $S(z)$ is complex analytic at all $\abs{z} < \rho$ hence by Proposition \ref{banach:ComplexAnalyticImpliesHolomorphic} we deduce that $S(z)$ is
  holmorphic at all $\abs{z} < \rho$.   In particular, all polynomials are entire.
\end{examp}

It turns out every holomorphic function is also complex analytic.   Proving this is one of our major goals, but getting to that result takes some time and requires some significant machinery involving curvilinear integrals.

For the next result note that any complex Banach space $X$ is also a
real Banach space, the metric on $\complexes$ given by absolute value
is the same as the Euclidean metric on $\reals^2$ (hence they have the
same topology) and therefore for any open $U \subset \complexes$ and $f : U
\to X$ we may also consider $f$ a function between real Banach spaces.

\begin{prop}\label{banach:HolomorphicFrechet} Let $X$ be a complex Banach space, $U \subset
  \complexes$ be open and $f : U \to X$ be a function.  Then $f$ is
  holomorphic at $z \in U$ if and only if $f$ thought of as a function
  between real Banach spaces is Frechet differentiable at $z$ and the
  derivative $Df(z) : \reals^2 \to X$ is $\complexes$-linear.   In this case, $f^\prime(z) = Df(z) \cdot \begin{pmatrix} 1 \\ 0 \\ \end{pmatrix}$.
\end{prop}
\begin{proof}
  First suppose that $f$ is holomorphic at $z$.   We claim that for
  $\begin{pmatrix} u \\ v \end{pmatrix} \in \reals^2$,
  \begin{align*}
    L \begin{pmatrix} u \\ v \end{pmatrix} &= \left(u + iv\right)
                                                 f^\prime(z)
  \end{align*}
  is the Frechet derivative.  Notationally it is much more convenient
  to just write an element of $\reals^2$ as a complex variable $h$ and then we are
  claiming that $L h = h f^\prime(z)$.   This follows since
  \begin{align*}
    \lim_{h \to 0} \norm{\frac{f(z + h) - f(z) - Lh}{\abs{h}}}
    &= \lim_{h \to 0} \norm{\frac{f(z + h) - f(z) - hf^\prime(z)}{h}} \\
    &= \lim_{h \to 0} \norm{\frac{f(z + h) - f(z)}{h} - f^\prime(z)} = 0
  \end{align*}

We continue to write elements of $\reals^2$ as complex numbers.   To see the converse suppose that $Df(z)$ is the $\complexes$-linear Frechet derivative
then we claim that $f^\prime(z) = Df(z) \cdot 1$.   By the assumption of $\complexes$-linearity we know that $h f^\prime(z) = Df(z) \cdot h$ for all $h$.   Now we just
reverse the previous computation to see that $f$ is holomorphic at $z$,
 \begin{align*}
   \lim_{h \to 0} \norm{\frac{f(z + h) - f(z)}{h} - f^\prime(z)} 
    &= \lim_{h \to 0} \norm{\frac{f(z + h) - f(z) - hf^\prime(z)}{h}} \\
    &=\lim_{h \to 0} \norm{\frac{f(z + h) - f(z) - Df(z) h}{\abs{h}}} = 0
  \end{align*}
\end{proof}

\section{Contour Integrals}

TODO: What assumptions do we make on $f$?   I.e. what restrictions on $f$ are needed so that $f(\gamma(t)) \gamma^\prime(t) $ is regulated?   It appears that we only need $f$ continuous (in which case $f(\gamma(t)) \gamma^\prime(t) $ is piecewise continuous hence regulated  by Proposition \ref{banach:PiecewiseContinuousIsRegulated}).

We need a notion of line integral for complex Banach space valued functions.   We base this definition on the Riemann integral in Banach spaces.  
\begin{defn}Given an open set $U \subset \complexes$, a continuous map $\gamma : [a,b] \to U$ is said to be \emph{piecewise differentiable} if there exist $a = a_0 < a_1 < \cdots < a_n = b$ such that $\gamma$ is continuously differentiable when restricted to each interval $(a_i, a_i+1)$ for $i=0, \dots, n-1$ and $\lim_{t \to a_i^+} \gamma^\prime(t)$ and $\lim_{t \to a_{i-1}^-} \gamma^\prime(t)$ exist.  Given a complex Banach space $X$ and a continuous function $f : U \to X$ define the \emph{contour integral} $\oint_\gamma f$ by the formula
\begin{align*}
\oint_\gamma f &= \sum_{i=0}^{n-1} \int_{a_i}^{a_{i+1}} f(\gamma(t)) \gamma^\prime(t) \, dt
\end{align*}
\end{defn}

The notation we have selected is potentially a little ambiguous since the definition might depend on the partition of the interval $[a,b]$.   It is easy to dispense with this concern.
\begin{prop}The value of the integral $\oint_\gamma f$ doesn't depend on the partition of $[a,b]$.
\end{prop} 
\begin{proof}
  Suppose $\gamma$ is piecewise differentiable with respect to the partition $a = a_0 < a_1 < \cdots < a_n = b$.  Now let $a \leq c \leq b$ and we consider adding the value $c$ to the existing partition.   Clearly if $a_i=c$ for $0 \leq i \leq n-1$, then the partition doesn't change.   Otherwise there is an $0 \leq i \leq n-1$ such that $a_i < c < a_{i+1}$.   By assumption $\gamma$ is continuously differentiable on each of $(a_i, c)$ and $(c, a_{i+1})$.   Moreover we also know that $\gamma^\prime(c)$ exists and $\gamma^\prime$ is continuous at $c$ hence
  \begin{align*}
    \lim_{t \to c^+} \gamma^\prime(t) = \gamma^\prime(c) = \lim_{t \to c^-} \gamma^\prime(t)
  \end{align*}
and it follows that $\gamma$ is piecewise differentiable with respect to $a = a_0 < \cdots < a_i < c < a_{i+1} < \cdots < a_n = b$.
  By Proposition \ref{banach:RiemannIntegralProperties}
\begin{align*}
\int_{a_i}^{a_{i+1}} f(\gamma(t)) \gamma^\prime(t) \, dt &= \int_{a_i}^c f(\gamma(t)) \gamma^\prime(t) \, dt + \int_c^{a_{i+1}} f(\gamma(t)) \gamma^\prime(t) \, dt
\end{align*}
so the value of $\oint_\gamma f$ is unchanged by the addition of the point $c$.   Suppose are given two different partitions $a = a_0 < a_1 < \cdots < a_n = b$ and $a = b_0 < b_1 < \cdots < b_m = b$ with $\gamma$ piecewise differentiable with respect to both.   Add the elements $b_0, b_1, \dots, b_m$ to the partition $a = a_0 < a_1 < \cdots < a_n = b$ and the integral remain unchanged.   Symmetrically we can start with $a = b_0 < b_1 < \cdots < b_m = b$ and add the $a_i$ without changing the integral.   At the end of both constructions, we have the integral with the partition defined by the union of the $a_0, a_1, \dots a_n$ and the $b_0, b_1, \dots b_m$, implying the two original integrals are equal.
\end{proof}

It turns out the integral only depends on the image of the curve $\gamma$. 
\begin{defn}Let $\gamma : [a,b] \to U$ and $\tilde{\gamma} : [\tilde{a}, \tilde{b}] \to U$ be piecewise differentiable curves then we say that $\tilde{\gamma}$ is a \emph{reparameterization} of $\gamma$ if there exists a continuously differentiable $\alpha : [a,b] \to [\tilde{a}, \tilde{b}]$ such that $\alpha^\prime(t) > 0$, $\alpha(a) = \tilde{a}$, $\alpha(b) = \tilde{b}$ and $\gamma(t) = \tilde{\gamma}(\alpha(t))$ for all $a \leq t \leq b$.
\end{defn}

\begin{prop}\label{banach:ContourIntegralParameterizationIndependence}Let $U \subset \complexes$ be open, $\gamma : [a,b] \to U$ be a piecewise differentiable curve and $\tilde{\gamma} : [\tilde{a}, \tilde{b}] \to U$ be a reparameterization of $\gamma$ then for any complex Banach space $X$ and continuous $f : U \to X$ we have 
\begin{align*}
\oint_{\gamma} f &= \oint_{\tilde{\gamma}} f
\end{align*}
\end{prop}
\begin{proof}
By the definition of the contour integral and restriction to an appropriate interval, it clearly suffices to consider the case in which $\gamma$ is assumed to be continuously differentiable so by the definition of the integral and the Chain Rule,
\begin{align*}
\oint_\gamma f &= \int_a^b f(\gamma(t)) \gamma^\prime(t) \, dt = \int_a^b f(\tilde{\gamma}(\alpha(t) ))  \tilde{\gamma}^\prime(\alpha(t)) \alpha^\prime(t) \, dt
\end{align*}
Now we perform a change of variables $s = \alpha(t)$ and by the Change of Variables formula (Theorem \ref{banach:RiemannIntegralChangeOfVariables}) we get
\begin{align*}
\oint_\gamma f &= \int_{\tilde{a}}^{\tilde{b}} f(\tilde{\gamma}(s)) \tilde{\gamma}^\prime(s) \, ds = \oint_{\tilde{\gamma}} f
\end{align*}
\end{proof}

\begin{prop}\label{banach:ContourIntegralOperations}Let $U \subset \complexes$ be open, $\gamma : [a,b] \to U$, $\tilde{\gamma} : [b, c] \to U$ with $\gamma(b) = \tilde{\gamma(b)}$.   Let $-\gamma : [a,b] \to U$ be defined by $-\gamma(t) = \gamma(b+a-t)$ and $\gamma + \tilde{\gamma} : [a,c] \to U$ be defined by
\begin{align*}
\left(\gamma + \tilde{\gamma}  \right)(t) = \begin{cases}
\gamma(t) & \text{for $a \leq t \leq b$} \\
\tilde{\gamma}(t) & \text{for $b < t \leq c$} \\
\end{cases}
\end{align*}
then each of $-\gamma$ and $\gamma + \tilde{\gamma}$ are piecewise differentiable and for all complex Banach spaces and $X$ and continuous $f : U \to X$
\begin{itemize}
\item[(i)] $\oint_{\gamma} f = - \oint_{-\gamma} f$
\item[(ii)] $\oint_{\gamma + \tilde{\gamma}} f = \oint_{\gamma} f + \oint_{\tilde{\gamma}} f$
\end{itemize}
\end{prop}
\begin{proof}
First observe that if $a=a_0 < a_1 < \dotsc < a_n=b$ is a partition for which $\gamma$ is continuously differentiable on $(a_i, a_{i+1})$ for all $0 \leq i < n$ then by the Chain Rule $-\gamma$ is also differentiable and therefore $-\gamma$ is piecewise differentiable.   By restriction to the appropriate interval we may assume that $\gamma$ and $\tilde{\gamma}$ are continuously differentiable on $[a,b]$.   Now just compute via the Chain Rule and the change of variables $s = a + b - t$,
\begin{align*}
\oint_{-\gamma} f &= \int_a^b f(\gamma(a + b -t)) \cdot -\gamma^\prime(a + b -t) \, dt = \int_b^a f(\gamma(s)) -\gamma^\prime(s) \, -ds \\
&= -\int_a^b f(\gamma(s)) \gamma^\prime(s) \, ds 
\end{align*}

Assume that $\gamma$ is piecewise differentiable with respect to partition $a=a_0 < a_1 < \cdots < a_n =b $ and $\tilde{\gamma}$ is piecewise differentiable with respect to the partition $b=a_n < a_{n+1} < \cdots < a_m$ with $m > n$.   Because $\gamma(b) = \tilde{\gamma}(b)$ it follows that $\gamma + \tilde{\gamma}$ is continuous.   It is also clear that $\gamma + \tilde{\gamma}$ is piecewise differentiable with respect to $a=a_0 < a_1 < \cdots < a_m = c$ because $\left(\gamma + \tilde{\gamma} \right)^\prime = \gamma^\prime$ on $(a_i, a_{i+1})$ for $0 \leq i < n$ and
$\left(\gamma + \tilde{\gamma} \right)^\prime = \tilde{\gamma}^\prime$ on $(a_i, a_{i+1})$ for $n \leq i < m$.   Now by Proposition \ref{banach:RiemannIntegralProperties} we see that
\begin{align*}
\oint_{\gamma + \tilde{\gamma}} f &= \int_a^c f\left(\left(\gamma + \tilde{\gamma}\right)(t)\right) \left(\gamma + \tilde{\gamma} \right)^\prime(t) \, dt \\
&=\int_a^b f\left(\left(\gamma + \tilde{\gamma}\right)(t)\right) \left(\gamma + \tilde{\gamma} \right)^\prime(t) \, dt  + \int_b^c f\left(\left(\gamma + \tilde{\gamma}\right)(t)\right) \left(\gamma + \tilde{\gamma} \right)^\prime(t) \, dt \\
&=\int_a^b f\left(\gamma(t)\right) \gamma^\prime(t) \, dt + \int_b^c f\left(\tilde{\gamma}(t)\right) \tilde{\gamma}^\prime(t) \, dt = \oint_{\gamma} f + \oint_{\tilde{\gamma}} f
\end{align*}
\end{proof}

\begin{prop}\label{banach:FundamentalTheoremOfCalculusForContourIntegrals}Let $U \subset \complexes$ be open, $\gamma : [a,b] \to U$, $X$ a complex Banach space and $f : U \to X$ holomorphic on $U$ then 
\begin{align*}
\oint_\gamma f^\prime &= f(\gamma(b)) - f(\gamma(a))
\end{align*}
\end{prop}
\begin{proof}
We apply the Chain Rule, Proposition \ref{banach:HolomorphicFrechet} and the Fundamental Theorem of
Calculus (Theorem \ref{FundamentalTheoremOfCalculusForBanachSpaceRiemannIntegrals}) 
\begin{align*}
\oint_\gamma f^\prime &= \int_a^b f^\prime(\gamma(t)) \gamma^\prime(t) \, dt = \int_a^b Df(\gamma(t)) \, dt =  f(\gamma(b)) - f(\gamma(a))
\end{align*}
\end{proof}



\begin{defn}Let $U \subset \complexes$ be open and $\gamma : [a,b] \to U$ be piecewise differentiable then we denote the \emph{arclength} of $\gamma$ by
  \begin{align*}
    \abs{\gamma} &= \int_a^b \abs{\gamma^\prime(t)} \, dt
  \end{align*}
\end{defn}

\begin{prop}\label{banach:ArclengthProperties}Let $U \subset \complexes$ be open and $\gamma : [a,b] \to U$ and $\tilde{\gamma} : [b,c] \to U$ be piecewise differentiable then
  \begin{itemize}
  \item[(i)] $\abs{\gamma} = \abs{-\gamma}$
  \item[(ii)] $\abs{\gamma + \tilde{\gamma}} = \abs{\gamma} + \abs{\tilde{\gamma}}$
  \item[(iii)] $\norm{\oint_{\gamma} f} \leq \sup_{a \leq t \leq b} \norm{f(\gamma(t))} \abs{\gamma} \leq \norm{f}_\infty \abs{\gamma}$ for all complex Banach spaces $X$ and bounded continuous $f : U \to X$.
  \end{itemize}
\end{prop}
\begin{proof}
  To see (i), use the change of variables $s = b+a-t$
  \begin{align*}
    \abs{\gamma} &= \int_a^b \abs{\gamma^\prime(t)} \, dt = -\int_b^a \abs{\gamma^\prime(b + a - s)} \, ds = \int_a^b \abs{\left(-\gamma\right)^\prime(s)} \, ds  = \abs{-\gamma}
  \end{align*}
  To see (ii), just use the definition of $\gamma + \tilde{\gamma}$ and linearity of the integral (Proposition \ref{banach:RiemannIntegralProperties}) to see
  \begin{align*}
    \abs{\gamma + \tilde{\gamma}} &= \int_a^c \abs{\left(\gamma + \tilde{\gamma}\right)^\prime(t)} \, dt \\
    &= \int_a^b \abs{\gamma^\prime(t)} \, dt + \int_b^c \abs{\tilde{\gamma}^\prime(t)} \, dt = \abs{\gamma} + \abs{\tilde{\gamma}}
  \end{align*}
  To see (iii), let $f : U \to X$ be given then use the Proposition \ref{NormRiemannIntegralBanachSpace}
  \begin{align*}
    \norm{\oint_\gamma f} &= \norm{\int_a^b f(\gamma(t)) \gamma^\prime(t) \, dt} \leq \int_a^b \norm{ f(\gamma(t)) \gamma^\prime(t)}  \, dt \\
    &\leq \sup_{a \leq t \leq b} \norm{f(\gamma(t))}  \int_a^b \abs{\gamma^\prime(t)} \, dt = \sup_{a \leq t \leq b} \norm{f(\gamma(t))}  \abs{\gamma}
  \end{align*}
  The fact that $\sup_{a \leq t \leq b} \norm{f(\gamma(t))} \leq \norm{f}_\infty = \sup_{z \in U} \norm{f(z)}$ follows by noting that $\gamma([a,b]) \subset U$.
\end{proof}

\begin{examp}\label{banach:ArclengthOfLineSegment}Let $z, w \in \complexes$ and let $\gamma : [a,b] \to \complexes$ be defined by $\gamma(t) = \frac{t-a}{b-a}w + \frac{b -t}{b-a}z$ be the
  line segment connecting $z$ and $w$ then $\abs{\gamma} = \abs{w -z}$ is the length of the line segment in the Euclidean metric.
\end{examp}
\begin{proof}
  The proof is an elementary computation
  \begin{align*}
    \abs{\gamma} &=  \int_a^b \abs{\gamma^\prime(t)} \, dt = \int_a^b \abs{\frac{w}{b-a} - \frac{z}{b-a}} \, dt = \abs{w - z}
   \end{align*}
\end{proof}

\section{Cauchy's Theorem and Cauchy's Integral Formula}

The most important tool in understanding holomorphic functions is the study of contour integrals of such a function.   In particular, the study of integrals of holomorphic functions on closed curves turns out be very powerful.   As we set out to give the basic theorems in this vein, it is convenient to start with the study of integrals over the boundary of rectangles.   In particular, we study rectangles whose sides are parallel to the real and imaginary axes of the complex plane.

\begin{defn}A \emph{rectangle} in the complex plane $\complexes$ is defined two complex numbers $z = a+bi$ and $w = c + di$ which satisfy $a < c$ and $b < d$.   Up to reparameterization the rectangle is defined by $\gamma : [0,4] \to \complexes$ with the formula
\begin{align*}
\gamma(t) &= \begin{cases}
(a + (c - a) t) + bi & \text{for $0 \leq t \leq 1$} \\
c + (b + (d-b)(t-1))i & \text{for $1 \leq t \leq 2$} \\
(c + (a - c) (t - 2)) + di & \text{for $2 \leq t \leq 3$} \\
a + (d + (b-d)(t-3))i & \text{for $3 \leq t \leq 4$} \\
\end{cases}
\end{align*}
If $U \subset \complexes$ is open then we that the rectangle is \emph{contained} in $U$ is $\gamma([0,4]) \subset U$.   We say that the rectangle is \emph{entirely contained} in $U$ if 
\begin{align*}
R &=[a,c] \times [b,d] = \lbrace u + vi \in U | a \leq u \leq c \text{ and } b \leq v \leq d \rbrace \subset U
\end{align*}
\end{defn}

Note that by Example \ref{banach:ArclengthOfLineSegment} and Proposition \ref{banach:ArclengthProperties} the arclength of the rectangle is equal to its perimeter in the middle school geometry sense.

\begin{tikzpicture}
[decoration={markings, 
    mark= at position 0.5 with {\arrow{stealth}}}
] 
\draw [postaction={decorate}] (0,0) -- (10,0);
\draw [postaction={decorate}] (10,0) -- (10,5);
\draw [postaction={decorate}] (10,5) -- (0,5);
\draw [postaction={decorate}] (0,5) -- (0,0);
\end{tikzpicture}

\begin{lem}(Goursat's Lemma)\label{banach:GoursatLemma}Let $X$ be a complex Banach space, $U \subset \complexes$ be an open set, and $f : U \to X$ be holomorphic at every $z \in U$, then for every rectangle $\gamma$ contained entirely in $U$ we have $\oint_\gamma f = 0$.   
\end{lem}
\begin{proof}
Let $R = [a,c] \times [b,d]$ as in the definition of rectangles.   Let $P=2\left( c + d - a -b \right)$ be the perimeter/arclength of the rectangle and let $D = \sqrt{(c-a)^2 + (d - b)^2}$ be the diameter of the rectangle (so that if $z,w \in R$ we have $\abs{z -w} \leq D$).

The main idea of the proof is to write the rectangle as a union of the four equally sized rectangles $R_1, R_2, R_3, R_4$ obtained by bisecting its sides.

\begin{tikzpicture}
[decoration={markings, 
    mark= at position 0.25 with {\arrow{stealth}},
    mark= at position 0.75 with {\arrow{stealth}}}
] 
\draw [postaction={decorate}] (0,0) -- (10,0);
\draw [postaction={decorate}] (10,0) -- (10,5);
\draw [postaction={decorate}] (10,5) -- (0,5);
\draw [postaction={decorate}] (0,5) -- (0,0);
\draw (5,0) -- (5,5);
\draw (0,2.5) -- (10,2.5);
\draw [->] (3.0,1.25) arc [start angle = 0, end angle = 270, radius = 0.5cm];
\draw [->] (8.0,1.25) arc [start angle = 0, end angle = 270, radius = 0.5cm];
\draw [->] (3.0,3.75) arc [start angle = 0, end angle = 270, radius = 0.5cm];
\draw [->] (8.0,3.75) arc [start angle = 0, end angle = 270, radius = 0.5cm];
\node at (7.5, 1.25) {$R_1$};
\node at (7.5, 3.75) {$R_2$};
\node at (2.5, 3.75) {$R_3$};
\node at (2.5, 1.25) {$R_4$};
\end{tikzpicture}

Note that each of the $R_i$ has perimeter $P/2$ and diameter $D/2$.   Note also if we let $\gamma_1, \gamma_2, \gamma_3, \gamma_4$ be boundaries their respective rectangles then by Proposition \ref{banach:ContourIntegralOperations}
\begin{align*}
\oint_\gamma f &= \oint_{\gamma_1} f + \oint_{\gamma_2} f + \oint_{\gamma_3} f + \oint_{\gamma_4} f
\end{align*}
(this follows because each interior edge is followed twice : once with each orientation).   By the triangle inequality, 
\begin{align*}
\norm{\oint_\gamma f} &\leq \norm{\oint_{\gamma_1} f} + \norm{\oint_{\gamma_2} f} + \norm{\oint_{\gamma_3} f} + \norm{\oint_{\gamma_4} f}
\end{align*}
and therefore for some $i$, we have $\norm{\oint_{\gamma_i} f} \geq \frac{1}{4} \norm{\oint_\gamma f}$.   

The labels on our subrectangles are arbitrary so let us assume $R_1$ is such a subrectangle and then inductively apply the same bisection procedure to yield a sequence of nested closed rectangles $R_1, R_2, \dots$ with boundaries $\gamma_1, \gamma_2, \dots$ such that for all $n \in \naturals$, the perimeter of $R_n$ is equal to $P/2^n$, the diameter of $R_n$ is $D/2^n$ and 
\begin{align*}
\norm{\oint_{\gamma_n} f} &\geq \frac{1}{4^n} \norm{\oint_{\gamma} f}
\end{align*}

By Lemma \ref{IntersectionCompactSets} we know that $\cap_{n=1}^\infty R_n \neq \emptyset$, so we may pick $z_0 \in \cap_{n=1}^\infty R_n$.  
We use the holomorphicity of $f$ at $z_0$ to bound the integral $\norm{\oint_{\gamma_n} f}$.   
Note that by the fact that polynomials are entire and Proposition \ref{banach:FundamentalTheoremOfCalculusForContourIntegrals} 
we know that $\oint_{\gamma_n} z \, dz=0$ and $\oint_{\gamma_n}dz = 0$ for all $n \in \naturals$.   
Using the fact that $f$ is holomorphic at $z_0$, let  $\epsilon > 0$ be given and pick $\delta >0$ such that for all $\abs{h} < \delta$,
\begin{align*}
\norm{f(z_0+h) - f(z_0) - h f^\prime(z_0)} &< \frac{\epsilon}{PD} \abs{h}
\end{align*}
Then pick $n$ large enough so that $D < 2^n \delta$, it follows from the norm bound for contour integrals in Proposition \ref{banach:ArclengthProperties}
\begin{align*}
\norm{\oint_{\gamma_n} f(z) \, dz} &= \norm{\oint_{\gamma_n}\left( f(z_0+z - z_0) - f(z_0) - (z - z_0) f^\prime(z_0)\right) \, dz }\\
&\leq \norm{f(z_0+z - z_0) - f(z_0) - (z - z_0) f^\prime(z_0)} \cdot \frac{P}{2^n}  \\
&< \frac{\epsilon}{PD} \abs{z - z_0} \cdot \frac{P}{2^n}  \leq \frac{\epsilon}{PD} \frac{D}{2^n} \frac{P}{2^n} = \frac{\epsilon}{4^n} \\
\end{align*}
Thus for $D < 2^n \delta$,
\begin{align*}
\norm{\oint_{\gamma} f} &\leq 4^n \norm{\oint_{\gamma_n} f} \leq \epsilon 
\end{align*}
Since $\epsilon>0$ was arbitrary, the result is proved.
\end{proof}

\begin{tikzpicture}
[decoration={markings, 
    mark= at position 0.25 with {\arrow{stealth}},
    mark= at position 0.75 with {\arrow{stealth}}}
] 
\draw (0,0) -- (10,0);
\draw (10,0) -- (10,5);
\draw (10,5) -- (0,5);
\draw (0,5) -- (0,0);
\draw (3,0) -- (3,5);
\draw (4,0) -- (4,5);
\draw (0,2) -- (10,2);
\draw (0,3) -- (10,3);
\node [circle, fill, inner sep=1pt, label=below right:$z_0$] at (3.5, 2.5) {};
\end{tikzpicture}

The following corollary of Goursat's Lemma appears to be an extension of sorts with relaxed hypotheses.   As we shall later see the seeming relaxation in the hypotheses is illusory and are, in fact, equivalent to the hypotheses in the original version of the lemma.
\begin{cor}\label{banach:GoursatLemmaPunctured}Let $X$ be a complex Banach space, $U \subset \complexes$ be an open set, $z_0 \in U$, and $f : U \to X$ be holomorphic at every $z \in U \setminus \lbrace z_0 \rbrace$ such that $\lim_{z \to z_0} (z - z_0) f(z) = 0$, then for every rectangle $\gamma$ contained entirely in $U$ with $z_0$ contained in the interior of $\gamma$, we have $\oint_\gamma f = 0$.   
\end{cor}
\begin{proof}
  Let $\epsilon>0$ be given and pick a $\delta>0$ such that $\norm{(z-z_0) f(z)} < \epsilon$ for all $\abs{z-z_0} < \frac{\delta}{\sqrt{2}}$.   It then follows that for the square $R_\epsilon$ of size $\delta$ and centered at $z_0$ we have $\norm{(z-z_0) f(z)} < \epsilon$ for all $z \in R_\epsilon$.   Let $\gamma_\epsilon$ be the perimeter of $R_\epsilon$ and note that by Goursat's Lemma and the holomorphicity of $f$ away from $z_0$ we have $\oint_\gamma f = \oint_{\gamma_\epsilon} f$ (refer to the picture above).   For all $z \in \gamma_\epsilon$ we have $\abs{z-z_0} \leq \delta/2$ and therefore
  \begin{align*}
    \norm{f(z)} &\leq \frac{\epsilon}{\abs{z - z_0}} \leq \frac{\epsilon}{\delta/2}
  \end{align*}
  Thus we have
  \begin{align*}
    \norm{\oint_{\gamma_\epsilon} f} &\leq \frac{\epsilon}{\delta/2} \abs{\gamma_\epsilon} = 8 \epsilon
  \end{align*}
Since $\epsilon>0$ was arbitrary we have $\oint_{\gamma_\epsilon} f = 0$.
\end{proof}

The previous results allow one to construct (local) antiderivatives (also called primitives) for holomorphic functions.
\begin{lem}\label{banach:LocalHolomorphicPrimitive}Let $X$ be a complex Banach space, $U \subset \complexes$ be an open set, and let $f: U \to X$ be continuous such that $\oint_\gamma f = 0$ for every rectangle $\gamma$ contained entirely in $U$.   Then for every open disk $D \subset U$ there exists a function $F$ holomorphic on $D$ such that $F^\prime = f$ on $D$.   
\end{lem}
\begin{proof}
  Let $z_0 = x_0 + iy_0 \in D$ be the center of the disk.  For every $z = x + i y\in D$ define the curve $\gamma_z : [0,2] \to D$ by
  \begin{align*}
    \gamma_z(t) &= \begin{cases}
      x_0 + i((1-t)y_0 + ty) & \text{for $0 \leq t \leq 1$} \\
      ((2-t)x_0 + (t-1) x) + iy & \text{for $1 \leq t \leq 2$} \\
    \end{cases}
  \end{align*}
as well as the curve $\tilde{\gamma}_z : [0,2] \to D$
  \begin{align*}
    \tilde{\gamma}_z(t) &= \begin{cases}
      ((1-t)x_0 + tx) + iy+0 & \text{for $0 \leq t \leq 1$} \\
      x + i((2-t)y_0 + (t-1) y) & \text{for $1 \leq t \leq 2$} \\
    \end{cases}
  \end{align*}
  and note that, in fact, the image of each curve does lie entirely in $D$.  TODO: Show this analytically; it is geometrically obvious.   By assumption, $\gamma_z - \tilde{\gamma}_z$ is a rectangle contained entirely in $U$ hence $\oint_{\gamma_z - \tilde{\gamma}_z} f = \oint_{\gamma_z} f  - \oint_{\tilde{\gamma}_z} f = 0$.    We define $F(z) = \oint_{\gamma_z} f  = \oint_{\tilde{\gamma}_z} f$.

  To see that $F$ is a primitive of $f$ we compute partial derivatives.   Let $h>0$ such that $z+h \in D$ and note that we have
  \begin{align*}
    F(z+h) - F(z) &= \oint_{\gamma_{z+h}} f -  \oint_{\gamma_z} f = \int_0^h f(z + t) \, dt
  \end{align*}
  Let $\epsilon>0$ be given, then since $f$ is continuous at $z$ we may pick $\delta>0$ such that $\norm{f(z+t) - f(z)} < \epsilon$ for all $\abs{t} < \delta$.   It follows that for $0 < h < \delta$,
  \begin{align*}
    \norm{\frac{\int_0^h f(z + t) \, dt }{h} - f(z)} &= \norm{\frac{\int_0^h \left( f(z + t) - f(z) \right) \, dt}{h}} \\
    &\leq \frac{\int_0^h \norm{ f(z + t) - f(z) } \, dt}{h} \leq \epsilon
  \end{align*}
  and therefore $\lim_{h \to 0^+} \frac{F(z+h) - F(z)}{h} = f(z)$.    For $h < 0$ we have $F(z) - F(z+h) = \int_h^0 f(z+t) \, dt$, letting $\epsilon$ and $\delta$ be as above, it follows for $-\delta < h < 0$,
 \begin{align*}
    \norm{\frac{\int_h^0 f(z + t) \, dt }{-h} - f(z)} &= \norm{\frac{\int_h^0 \left( f(z + t) - f(z) \right) \, dt}{-h}} \\
    &\leq \frac{\int_h^0 \norm{ f(z + t) - f(z) } \, dt}{-h} \leq \epsilon
  \end{align*}
  and therefore $\lim_{h \to 0^-} \frac{F(z+h) - F(z)}{h} = \lim_{h \to 0^-} \frac{\int_h^0 f(z + t) \, dt }{-h}  = f(z)$.   This shows that $D_1 F(z) = f(z)$ hence is continuous.

Arguing analogously using $\tilde{\gamma}_z$ we see that for $h>0$ such that $z + ih \in D$,
  \begin{align*}
    F(z+ih) - F(z) &= \oint_{\tilde{\gamma}_{z+ih}} f -  \oint_{\tilde{\gamma}_z} f = \int_0^h f(z + i t) i \, dt = i \int_0^h f(z + i t) \, dt
  \end{align*}
Letting $\epsilon$ and $\delta$ be as above we get 
 \begin{align*}
    \norm{\frac{\int_0^h f(z + it) \, dt }{h} - f(z)} &= \norm{\frac{\int_0^h \left( f(z + it) - f(z) \right) \, dt}{h}} \\
    &\leq \frac{\int_0^h \norm{ f(z +i t) - f(z) } \, dt}{h} \leq \epsilon
  \end{align*}
  hence $\lim_{h \to 0^+} \frac{F(z+ih) - F(z)}{h} = i f(z)$.   The argument for $h<0$ proceeds the same way and we conclude $D_2 F(z) = i f(z)$ hence is continuous.   By Proposition \ref{PartialDerivativesBanachSpaces} we know that $F$ is continuously Frechet differentiable and
  \begin{align*}
    DF(z) \cdot \begin{pmatrix} x \\ y \end{pmatrix} &= D_1F(z) \cdot x + D_2 F(z) \cdot y = f(z) x + i f(z) y
  \end{align*}
  which shows that $DF(z) \cdot w= w f(z)$ is complex linear.   By Proposition \ref{banach:HolomorphicFrechet} we see that $F$ is holomorphic and $F^\prime = f$.
\end{proof}

\begin{cor}\label{banach:PointContourIntegralZero}Let $X$ be a complex Banach space, $U \subset \complexes$ be an open set, and let $f: U \to X$ be continuous such that $\oint_\gamma f = 0$ for every rectangle $\gamma$ contained entirely in $U$.   For any point $z \in U$, $\oint_{\lbrace z \rbrace} f = 0$.
\end{cor}
\begin{proof}
  Parameterize the point contour $\lbrace z \rbrace$ by $\gamma : [0,1] \to U$ with $\gamma(t) = z$ for all $0 \leq t \leq 1$.   Apply Lemma \ref{banach:LocalHolomorphicPrimitive} to get an open disc $D \subset U$ containing $z$ and a holomorphic function $F : D \to X$ such that $F^\prime(w) = f(w)$ for all $w \in D$.   Then by Proposition \ref{banach:FundamentalTheoremOfCalculusForContourIntegrals}, we have
  \begin{align*}
     \oint_{\lbrace z \rbrace} f &= \oint_{\gamma} F^\prime = F(\gamma(1)) - F(\gamma(0)) = F(z) - F(z) = 0 
  \end{align*} 
\end{proof}

The next step in the development of Cauchy's Theorem is the extension of the result to more general curves.   For that purpose we introduce an important topological
idea.
\begin{defn}Let $X$ be a topological space and let $\gamma_0 : [a,b] \to X$ and $\gamma_1: [a,b] \to X$ be continuous with $\gamma_0(a) = \gamma_0(b)$ and $\gamma_1(a) = \gamma_1(b)$.   We say that $\gamma_0$ and $\gamma_1$ are \emph{homotopic} if there exists a continuous map $H : [a,b] \times [0,1] \to X$ such that
  \begin{itemize}
  \item[(i)] $H(a,t) = H(b,t)$ for all $0 \leq t \leq 1$
  \item[(ii)] $H(s,0) = \gamma_0(s)$ for all $a \leq s \leq b$
  \item[(iii)] $H(s,1) = \gamma_1(s)$ for all $a \leq s \leq b$
  \end{itemize}
\end{defn}

Our goal is to show that contour integrals of a holomorphic function over homotopic curves are equal.   To get there, first we need to use some topology to extend our construction of primitives from a purely local construct to one with a more global character.   The arguments that follow require the following sequence of lemmata.

\begin{lem}\label{banach:EuclideanRectangleThickening}Let $F$ be a closed set in $\reals^d$, let $R = [s_1, t_1] \times \dotsc \times [s_d, t_d] \subset F$ and for any $\epsilon > 0$ denote $R_\epsilon = [s_1-\epsilon, t_1+ \epsilon] \times \dotsc \times [s_d - \epsilon, t_d + \epsilon] \subset \reals^d$.   Let $U \subset F$ be open in the relative topology such that $R \subset U$,
  then there exists $\epsilon > 0$, such that $R_\epsilon \cap F \subset U$.
\end{lem}
\begin{proof}
  Since $U$ is relatively open in $F$ there exists an open set $V \subset \reals^d$ such that $U = V \cap F$.   Clearly $R \subset V$ and since $V$ is open we know that $V^c$ is closed and $d(x, V^c)$ defines a Lipschitz function (hence uniformly continuous) on $R$.   By the Heine-Borel Theorem \ref{HeineBorel} we know that $R$ is compact and then by Theorem \ref{ExistenceGlobalMinimizerContinuous} we know that there exists $x^* \in R$ such that $d(x^*, V^c) \leq d(x, V^c)$ for all $x \in R$.   Moreover since $V^c$ is closed and $R \cap V^c = \emptyset$ we know that $0 < d(x^*, V^c)$.   Let $\epsilon = d(x^*, V^c)/3\sqrt{d}$, then for any $y \in R_\epsilon$, we have $d(y, R) \leq \sqrt{d \epsilon^2} \leq \sqrt{d} \epsilon = d(x^*, V^c)/3$.
   Pick an $x \in R$ such that $d(y, x) \leq d(y, R) + d(x^*, V^c)/3 $.   By the triangle inequality, for every $x \in R$, and $z \in V^c$,
  \begin{align*}
    d(x^*, V^c) &\leq d(x, V^c) \leq d(x, z) \leq d(x, y) + d(y, z) \\
                &\leq d(y, R) + d(x^*, V^c)/3 + d(y, V^c) \\
                &\leq  \frac{2}{3}d(x^*, V^c) + d(y, V^c) 
  \end{align*}
  Thus $0 < d(x^*,V^c)/3 \leq d(y, V^c)$ therefore $y \in V$.    Therefore $R_\epsilon \subset V$ and $R_\epsilon \cap F \subset V \cap F = U$.
\end{proof}

\begin{lem}\label{banach:PrimitiveFollowingMappingSubdivision}Let $R= [a_1, b_1] \times \cdots \times [a_d, b_d]$ be a closed rectangle in $\reals^d$ and let $\lbrace U_\alpha \rbrace$ be an open covering of $R$.    There exists finite sequences $a_1 = s^1_0 < s^1_1 < \cdots < s^1_{n_1 + 1} = b_1, \dots, a_d = s^d_0 < s^d_1 < \cdots < s^d_{n_d + 1} = b_d$ such that for every $0 \leq i_j \leq n_j$ for $1\leq j \leq d$ the rectangle $[s^1_{i_1}, s^1_{i_i+1}] \times \cdots \times [s^d_{i_d}, s^d_{i_d+1}]$ is contained in some $U_\alpha$.   
\end{lem}
\begin{proof}
  Let $r = \sqrt{(b_1 - a_1)^2 + \cdots + (b_d-a_d)^2}$ so that $R$ is contained in the open ball of radius $r$ centered at the midpoint of $R$.   Dyadically subdivide the rectangle.   That is to say, for $n > 0$ subdivide each side $1 \leq j \leq d$ of the rectangle into $2^n$ equally sized parts of length $(b_j-a_j)/2^n$ : $s^j_{i,n} = a_j + i(b_j-a_j)/2^n$ for $0 \leq i \leq 2^n$.   Arguing by contradiction, suppose that for every $n > 0$ there exists a dyadic rectangle that is not contained in any of the $\lbrace U_\alpha \rbrace$.   Let that rectangle be $R_n = [s^1_{i_1,n}, s^1_{i_1+1,n}] \times \cdots \times [s^d_{i_d,n}, s^d_{i_d+1,n}]$; let $x_n$ be the center of $R_n$ and note that $R_n \subset B(x_n, r/2^n) \cap R$.   By the Heine-Borel Theorem \ref{HeineBorel} we know that $R$ is compact and therefore the exists a convergent subsequence of $x_n$.   Let $x \in R$ be the limit of the convergent subsequence.

  Let $\epsilon>0$ be given and pick $n$ sufficiently large such that $\abs{x_n - x} < \epsilon/2$ and $r < \epsilon 2^{n-1}$; then by the triangle inequality $B(x_n, r/2^n) \subset  B(x, \epsilon)$.   It follows that $R_n \subset B(x_n, r/2^n) \cap R \subset B(x, \epsilon) \cap R$ which shows that $B(x, \epsilon) \cap R$ is not contained in any $\lbrace U_\alpha \rbrace$.   Since $\epsilon>0$ was arbitrary this contradicts the fact the  $\lbrace U_\alpha \rbrace$ are an open subcover of $R$.
\end{proof}

\begin{lem}\label{banach:PrimitiveFollowingMappingThickening}Let $R= [a_1, b_1] \times \cdots \times [a_d, b_d]$ be a closed rectangle in $\reals^d$ and let $a_1 = s^1_0 < s^1_1 < \cdots < s^1_{n_1 + 1} = b_1, \dots, a_d = s^d_0 < s^d_1 < \cdots < s^d_{n_d + 1} = b_d$ be finite partitions such that for every $0 \leq i_j \leq n_j$ for $1\leq j \leq d$ the rectangle $R_{i_1, \dots, i_d} = [s^1_{i_1}, s^1_{i_i+1}] \times \cdots \times [s^d_{i_d}, s^d_{i_d+1}]$ is contained in a relatively open set $U_{i_1, \dots, i_d}$ in $R$.   There exists an $\epsilon>0$ such that the thickening $R_{i_1, \dots, i_d, \epsilon} = [s^1_{i_1}  - \epsilon, s^1_{i_i+1}+\epsilon] \times \cdots \times [s^d_{i_d} - \epsilon, s^d_{i_d+1} + \epsilon] \cap R$ is contained in $U_{i_1, \dots, i_d}$.   Morever $R_{i_1, \dots, i_d, \epsilon}  \cap R_{j_1, \dots, j_d, \epsilon} =\emptyset$ whenever there exists $1 \leq k \leq d$ such that $\abs{i_k, j_k} \geq 2$.
\end{lem}
\begin{proof}
  Let $\mathcal{I} = \lbrace (i_1, \dots, i_d) \mid 0 \leq i_j \leq n_j \rbrace$.   Now for each $I \in \mathcal{I}$ with $I = (i_1, \dots, i_d)$, let $R_I = [s^1_{i_1}, s^1_{i_i+1}] \times \cdots \times [s^d_{i_d}, s^d_{i_d+1}]$ be a rectangle and select a particular open neighborhood $U_\alpha$ and call it $U_I$ (so $R_I \subset U_I$).  Apply Lemma \ref{banach:EuclideanRectangleThickening} to each $R_I$ to get an $\epsilon_I > 0$ such that the thickening $R_{I, \epsilon_I} \subset U_I$.   Also let $0 < M = \min_{1 \leq j \leq d} \min_{0 \leq i \leq n_i} (s^j_{i+1} - s^j_i)$ be the minimum distance between adjacent partition element on any of the axes.  Let $\epsilon = \min_{I \in \mathcal{I}} \epsilon_I \wedge M/2$.   For each $I \in \mathcal{I}$, $R_{I, \epsilon} \subset R_{I, \epsilon_I} \subset U_I$.

  Suppose that $I, J \in \mathcal{I}$ and there exists $1 \leq k \leq d$ such that $i_k +2 \leq j_k$.   If $x \in R_{I, \epsilon_I}$ and $x \in R_{J, \epsilon_J}$ then
  \begin{align*}
    \norm{x - y} &\geq \abs{x_k - y_k} \geq \abs{(s_{j_k} - \epsilon) - (s_{i_k} + \epsilon)} \geq 2M - 2\epsilon \geq M > 0
  \end{align*}
  and therefore $R_{I, \epsilon_I} \cap R_{J, \epsilon_J} = \emptyset$.
\end{proof}

\begin{defn}Let $X$ be a complex Banach space, $U \subset \complexes$ be an open set, let $f: U \to X$ be continuous such that $\oint_\gamma f = 0$ for every rectangle $\gamma$ contained entirely in $U$ and suppose $\gamma : [a,b] \to U$ is continuous.   A \emph{holomorphic primitive along $\gamma$} is a continuous function $\lambda : [a,b] \to X$ such that for every $a \leq t \leq b$ there exists an open neighborhood $\gamma(t) \in U_t\subset \complexes$, a holomorphic function $F_t : U_t \to X$ and $\epsilon_t > 0$ such that 
  \begin{itemize}
  \item[(i)] $F_t^{\prime}(z) = f(z)$ for all  $z \in U_t$
  \item[(ii)] $F_t(\gamma(s)) = \lambda(s)$ for $s \in (t - \epsilon_t, t + \epsilon_t) \cap [a,b]$.
  \end{itemize}
\end{defn}

\begin{prop}\label{banach:HolomorphicPrimitive}Let $X$ be a complex Banach space, $U \subset \complexes$ be an open set, let $f: U \to X$ be continuous such that $\oint_\gamma f = 0$ for every rectangle $\gamma$ contained entirely in $U$ and suppose $\gamma : [a,b] \to U$ is continuous.   Then a holomorphic primitive $\lambda$ for $f$ along $\gamma$ exists and is unique up to the addition of a constant.
  Moreover, if $\gamma$ is piecewise smooth, we also know that
  \begin{align*}
    \oint_{\gamma} f &= \lambda(b) - \lambda(a)
  \end{align*}
\end{prop}
\begin{proof}
  By Lemma \ref{banach:LocalHolomorphicPrimitive} we know that for every $a \leq t \leq b$ there exists an open neighborhood $\gamma(t) \in U_t$ and a holomorphic function $\tilde{F}_t : U \to X$ such that $\tilde{F}^\prime_t = f$ on $U_t$.   By continuity of $\gamma$ we know that $\gamma^{-1}(U_t)$ is an open set containing $t$ (e.g. Lemma \ref{ContinuousPreImageOfOpenSet}).   Apply Lemma \ref{banach:PrimitiveFollowingMappingSubdivision} to the open cover $\gamma^{-1}(U_t)$ and get a sequence $a=t_0 < t_1 < \dotsc < t_{n+1} =b$ such that $[t_j, t_{j+1}] \subset \gamma^{-1}(U_t)$ for some $a \leq t \leq b$.   For each $0 \leq j \leq n$ pick such an open set and call it $U_j$ so that $\gamma([t_j, t_{j+1}] \subset U_{j}$.  Also let $\tilde{F}_j$ be the corresponding primitive of $f$ so that $\tilde{F}^\prime_j (z) = f(z)$ for every $z \in U_j$.   Now apply Lemma \ref{banach:PrimitiveFollowingMappingThickening} to get an $\epsilon > 0$ such that $(t_j - \epsilon, t_j + \epsilon) \cap [a,b] \subset \gamma^{-1}(U_j)$ for all $0 \leq j \leq n$.

  We now use the constructed $t_i$ to define $\lambda$ and $F_{t}$ by induction over the intervals $[t_j, t_{j+1}]$.
 We begin by defining
$F_{0} =  \tilde{F}_{a}$  and for $a=t_0 \leq t \leq t_1$ define $U_t = U_{0}$ and $F_t = F_0$.   For $a=t_0 \leq t < t_1 + \epsilon$ define
$\lambda(t) = F_0(\gamma(t))$ (this is a valid definition because $\gamma([t_0, t_1+\epsilon) ) \subset U_0$).  By construction for all
$F^\prime_{t} (z) = F^\prime_{0} (z) = f(z)$ for all $z \in U_t = U_0$.  Moreover, if $a=t_0 \leq t \leq t_1$ then $(t - \epsilon, t + \epsilon) \cap [a,b] \subset [t_0, t_1 + \epsilon)$
and it follows by definition of $\lambda$ that $F_t(\gamma(s)) = \lambda(s)$ for all $s \in (t - \epsilon, t + \epsilon) \cap [a,b]$.
Since by assumption $\gamma$ is continuous and, as a holomorphic function, $F_{0}$ is continuous it follows that $\lambda$ is continuous on $[t_0, t_1+\epsilon)$.

  For the induction step, assume that for $(s,t) \in R_{i,0} \cup \dotsc \cup R_{i, j-1} = [s_i, s_{i+1}] \times [t_0, t_j]$ we have defined $\overline{F}_{s,t}$  and $U_{s,t}$ such that (i) holds.  Also assume that $\lambda_i$ has been defined on $R_{i,0,\epsilon} \cup \dotsc \cup R_{i,j-1, \epsilon}$ such that (ii) holds for all $(s,t) \in R_{i,0} \cup \dotsc \cup R_{i, j-1} = [s_i, s_{i+1}] \times [t_0, t_j]$.   In addition we assume that there exist open sets $U_{i,0}, \dots, U_{i, j-1}$ and holomorphic functions $\overline{F}_{i,0}, \dots, \overline{F}_{i, j-1}$ such that $U^i_{s,t} = U_{i,k}$ and $\overline{F}_{s,t} = \overline{F}_{i, k}$ for $(s,t) \in [s_i, s_{i+1}] \times [t_k, t_{k+1}]$.   TODO: Disambiguate the definition at the points $t_k$.   Now, consider the rectangle $R_{i, j}$.  We know that $\gamma(R_{i,j}) \subset U_{i,j}$ and we have primitive $\tilde{F}_{i,j}$ on $U_{i,j}$.   We know that $\gamma(s_i, t_j) \in U_{i,j-1} \cap U_{i,j}$ and therefore $U_{i,j-1} \cap U_{i,j} \neq \emptyset$.   Furthermore we know from above that since $U_{i,j-1}$ and $U_{i,j}$ are disks hence $U_{i,j-1} \cap U_{i,j}$ is connected.   We know that $\overline{F}_{i, j-1}$ and $\tilde{F}_{i,j}$ are holomorphic on $U_{i,j-1} \cap U_{i,j}$ and $\left(\overline{F}_{i, j-1} - \tilde{F}_{i,j}\right)^\prime = f - f = 0$ hence by TODO: we know that there exists $c \in X$ such that $\overline{F}_{i, j-1} - \tilde{F}_{i,j} = c$ on $U_{i,j-1} \cap U_{i,j}$.   Define $\overline{F}_{i,j} = \tilde{F}_{i,j} + c$ on $U_{i,j}$ so that $\overline{F}_{i, j-1} = \overline{F}_{i,j}$ on $U_{i,j-1} \cap U_{i,j}$.   For $(s,t) \in  [s_i, s_{i+1}]  \times (t_j, t_{j+1}]$, define $U^i_{s,t} = U_{i,j}$ and $\overline{F}_{s,t} = \overline{F}_{i,j}$.   By construction (i) holds on $R_{i,0} \cup \dotsc \cup R_{i, j}$.   Now we need to extend the definition of $\lambda_i$.   On $R_{i,j,\epsilon}$, define $\lambda_i(s,t) = \overline{F}_{i,j}(\gamma(s,t)$.   Arguing as above, it is clear that as a composition of continuous functions, $\lambda_i$ is continuous on $R_{i,j,\epsilon}$.   The first issue we have to resolve is that this definition does not conflict with the previous definiton of $\lambda_i$.   In particular, by the induction hypothesis $\lambda_i(s,t) = \overline{F}_{i,j-1}(\gamma(s,t))$ on $R_{i,j-1,\epsilon}$, so we need to make sure that the new definition agrees with the old one on $R_{i,j-1,\epsilon} \cap R_{i,j,\epsilon}$.   This follows because $\gamma(R_{i,j-1,\epsilon}\cap R_{i,j,\epsilon}) \subset U_{i,j-1} \cap U_{i,j}$ and we have defined $\overline{F}_{i,j}$ so that it equals $\overline{F}_{i,j-1}$ on $U_{i,j-1} \cap U_{i,j}$.   If $j>1$, then we observe that by Lemma \ref{banach:PrimitiveFollowingMappingThickening} $R_{i,k,\epsilon}\cap R_{i,j,\epsilon} = \emptyset$ for all $k \leq j-2$ and so there is no possible conflict from those parts of the domain of $\lambda_i$.

Next we show the uniqueness of $\lambda$ up to an additive constant.   Suppose $\lambda_1$ and $\lambda_2$ are two continuous functions satisfying the conclusion of the result.   There exist open disks $U_1$ and $U_2$ of $\gamma(a)$ and holomorphic functions $F_1 : U_1 \to X$ and $F_2 : U_2 \to X$ such that $F_1^\prime = F_2^\prime = f$ on $U_1 \cap U_2$ and $F_1(\gamma(a)) = \lambda_1(a)$ and $F_2(\gamma(a)= \lambda_2(a)$.   Since we know that $F_1 - F_2$ is constant on $U_1 \cap U_2$, by redefining $\lambda_2$ and $F_2$ by adding a constant we may assume that $\lambda_1$ and $\lambda_2$ agree on a relatively open neighborhood of $a$.   We want to show that $\lambda_1$ is equal to $\lambda_2$ on the entirety of $[a,b]$.  

TODO: Move this to an exercise.
\begin{clm}Intersection of two disks in a normed vector space is connected.
\end{clm}
It suffices to show that a ball contains the line segment containing any two points in the ball.  This follows from the following computation which is valid in any normed vector space $0 \leq t \leq 1$,
\begin{align*}
  \norm{tz + (1-t) w - x} &\leq \norm{tz - tx + (1-t) w - (1-t) x} \\
  &\leq t \norm{z - x}  + (1-t)\norm{w - x}  \leq \norm{z - x} \vee \norm{w - x}
\end{align*}
\end{proof}

It is important to note that the existence of holomorphic primitives along a curve applies to general continuous (not necessarily piecewise smooth) $\gamma : [a,b] \to U$.
The uniqueness up to additive constant of holomorphic primitive $\lambda$ along a curve shows that that the value $\lambda(b) - \lambda(a)$ depends only on the underlying $\gamma$.   Lastly Proposition \ref{banach:HolomorphicPrimitive} shows that the following definition is consistent with our original definition of contour integrals for piecewise smooth $\gamma$.

\begin{defn}Let $X$ be a complex Banach space, $U \subset \complexes$ be an open set, let $f: U \to X$ be continuous such that $\oint_\gamma f = 0$ for every rectangle $\gamma$ contained entirely in $U$ and suppose $\gamma : [a,b] \to U$ is continuous.   We define the contour integral
  \begin{align*}
    \oint_{\gamma} f &= \lambda(b) - \lambda(a)
  \end{align*}
  where $\lambda$ is a holomorphic primitive along $\gamma$.
\end{defn}

\begin{prop}\label{banach:HolomorphicPrimitiveOnRectangle}Let $X$ be a complex Banach space, $U \subset \complexes$ be an open set, let $f: U \to X$ be continuous such that $\oint_\gamma f = 0$ for every rectangle $\gamma$ contained entirely in $U$ and suppose $\gamma : [a,b] \times [c,d] \to U$ is continuous.   There exists a continuous function $\lambda : [a,b] \times [c,d] \to X$ unique up to addition by a constant and an $\epsilon>0$ such that for every $a \leq s \leq b$,  $c \leq t \leq d$ there exists an open neighborhood $\gamma(s,t) \in U_{s,t}\subset \complexes$ and a holomorphic function $F_{s,t} : U_{s,t} \to X$ such that 
  \begin{itemize}
  \item[(i)] $F_{s,t}^{\prime}(z) = f(z)$ for all  $z \in U_{s,t}$
  \item[(ii)] $\gamma((s - \epsilon, s + \epsilon) \times (t - \epsilon,    t+\epsilon) \cap [a,b] \times [c,d]) \subset U_{s,t}$ 
  \item[(iii)] $F_{s,t}(\gamma(u,v)) = \lambda(s,t)$ for $u \in (s - \epsilon, s + \epsilon) \cap [a,b]$ and $v \in (t - \epsilon, t + \epsilon) \cap [c,d]$ .
  \end{itemize}
\end{prop}
\begin{proof}
  By Lemma \ref{banach:LocalHolomorphicPrimitive} we know that for every $(s,t) \in [a,b] \times [c,d]$ there exists an open neighborhood $\gamma(s,t) \in U_{s,t}$ and a holomorphic function $\tilde{F}_{s,t} : U_{s,t} \to X$ such that $\tilde{F}^\prime_{s,t} = f$ on $U_{s,t}$.   By continuity of $\gamma$ we know that $\gamma^{-1}(U_{s,t})$ is an open set containing $(s,t)$ (e.g. Lemma \ref{ContinuousPreImageOfOpenSet}).   Apply Lemma \ref{banach:PrimitiveFollowingMappingSubdivision} to the open cover $\gamma^{-1}(U_{s,t})$ and get a sequence $a=s_0 < s_1 < \dotsc < s_{n+1} =b$ and a sequence $c=t_0 < t_1 < \dotsc < t_{m+1} =d$ such that $R_{i,j} = [s_i, s_{i+1}] \times [t_j, t_{j+1}] \subset \gamma^{-1}(U_{s,t})$ for some $(s,t) \in [a,b] \times [c,d]$.  For simplicity, denote that set $U_{i,j}$ so that $\gamma([s_i, s_{i+1}] \times [t_j, t_{j+1}]) \subset U_{i,j}$.  Also denote the corresponding primitive as $\tilde{F}_{i,j}$.    Next apply Lemma \ref{banach:PrimitiveFollowingMappingThickening} to get an $\epsilon>0$ such that $R_{i,j, \epsilon} = [s_i-\epsilon, s_{i+1}+ \epsilon] \times [t_j - \epsilon, t_{j+1}+\epsilon] \cap [a,b] \times [c,d] \subset \gamma^{-1}(U_{i,j})$ for all $0 \leq i \leq n$ and $0 \leq j \leq m$.  

  \begin{clm}For each $0 \leq i \leq n$ there exist holomorphic functions $\overline{F}_{i,j} : U_{i,j} \to X$ and a continuous function $\lambda_i : R_{i,0,\epsilon} \cup \dotsc \cup R_{i,m,\epsilon} \to X$ such that if we define
    \begin{align*}
      \overline{F}_{s,t} &= \begin{cases}
        \overline{F}_{i,0} & \text{ for  $(s,t) \in [s_i, s_{i+1}] \times [t_0, t_1]$} \\
        \overline{F}_{i,j} & \text{ for $1 \leq j \leq m$ and  $(s,t) \in [s_i, s_{i+1}] \times (t_j, t_{j+1}]$} \\
        \end{cases} \\
      U^i_{s,t} &= \begin{cases}
        U_{i,0} & \text{ for  $(s,t) \in [s_i, s_{i+1}] \times [t_0, t_1]$} \\
        U_{i,j} & \text{ for $1 \leq j \leq m$ and  $(s,t) \in [s_i, s_{i+1}] \times (t_j, t_{j+1}]$} \\
      \end{cases}
    \end{align*}
then for all $(s,t) \in [s_i, s_{i+1}] \times [c,d]$,     
  \begin{itemize}
  \item[(i)] $\overline{F}_{s,t}^{\prime}(z) = f(z)$ for all  $z \in U^i_{s,t}$
  \item[(ii)] $\overline{F}_{s,t}(\gamma(u,v)) = \lambda_i(s,t)$ for $u \in (s - \epsilon, s + \epsilon) \cap [a,b]$ and $v \in (t - \epsilon, t + \epsilon) \cap [c,d]$ .
  \end{itemize}
 \end{clm}
  To prove the claim, we perform a construction analogous to that in Proposition \ref{banach:HolomorphicPrimitive} on each strip $[s_i, s_{i+1}] \times [c,d]$.   To begin define $\overline{F}_{i,0} = \tilde{F}_{i,0}$ so it is holomorphic on $U_{i,0}$ and $\overline{F}^\prime_{i,0} = f$ on $U_{i,0}$.   For $(s,t) \in R_{i,0}$, define $\overline{F}_{s,t} = \overline{F}_{i,0}$ and $U_{s,t} = U_{u,0}$.   As noted, (i) holds for such $(s,t)$.  For $(s,t) \in R_{i, 0, \epsilon}$, define $\lambda_i(s,t) = F_{i,0}(\gamma(s,t))$.   For $(s,t) \in R_{i,0}$ we have for any $(s,t) \in R_{i,0}$ and $u \in (s - \epsilon, s + \epsilon) \cap [a,b]$ and $v \in (t - \epsilon, t + \epsilon) \cap [c,d]$, $(u,v) \in R_{i,0, \epsilon}$ hence $]\lambda_i(u,v) = \overline{F}_{i,0}(\gamma(u,v))$ by definition.   Since by assumption $\gamma$ is continuous and as a holomorphic function $\overline{F}_{i,0}$ is continuous it follows that $\lambda_i$ is continuous on $R_{i,0, \epsilon}$.
  For the induction step, assume that for $(s,t) \in R_{i,0} \cup \dotsc \cup R_{i, j-1} = [s_i, s_{i+1}] \times [t_0, t_j]$ we have defined $\overline{F}_{s,t}$  and $U_{s,t}$ such that (i) holds.  Also assume that $\lambda_i$ has been defined on $R_{i,0,\epsilon} \cup \dotsc \cup R_{i,j-1, \epsilon}$ such that (ii) holds for all $(s,t) \in R_{i,0} \cup \dotsc \cup R_{i, j-1} = [s_i, s_{i+1}] \times [t_0, t_j]$.   In addition we assume that there exist open sets $U_{i,0}, \dots, U_{i, j-1}$ and holomorphic functions $\overline{F}_{i,0}, \dots, \overline{F}_{i, j-1}$ such that $U^i_{s,t} = U_{i,k}$ and $\overline{F}_{s,t} = \overline{F}_{i, k}$ for $(s,t) \in [s_i, s_{i+1}] \times [t_k, t_{k+1}]$.   TODO: Disambiguate the definition at the points $t_k$.   Now, consider the rectangle $R_{i, j}$.  We know that $\gamma(R_{i,j}) \subset U_{i,j}$ and we have primitive $\tilde{F}_{i,j}$ on $U_{i,j}$.   We know that $\gamma(s_i, t_j) \in U_{i,j-1} \cap U_{i,j}$ and therefore $U_{i,j-1} \cap U_{i,j} \neq \emptyset$.   Furthermore we know from above that since $U_{i,j-1}$ and $U_{i,j}$ are disks hence $U_{i,j-1} \cap U_{i,j}$ is connected.   We know that $\overline{F}_{i, j-1}$ and $\tilde{F}_{i,j}$ are holomorphic on $U_{i,j-1} \cap U_{i,j}$ and $\left(\overline{F}_{i, j-1} - \tilde{F}_{i,j}\right)^\prime = f - f = 0$ hence by TODO: we know that there exists $c \in X$ such that $\overline{F}_{i, j-1} - \tilde{F}_{i,j} = c$ on $U_{i,j-1} \cap U_{i,j}$.   Define $\overline{F}_{i,j} = \tilde{F}_{i,j} + c$ on $U_{i,j}$ so that $\overline{F}_{i, j-1} = \overline{F}_{i,j}$ on $U_{i,j-1} \cap U_{i,j}$.   For $(s,t) \in  [s_i, s_{i+1}]  \times (t_j, t_{j+1}]$, define $U^i_{s,t} = U_{i,j}$ and $\overline{F}_{s,t} = \overline{F}_{i,j}$.   By construction (i) holds on $R_{i,0} \cup \dotsc \cup R_{i, j}$.   Now we need to extend the definition of $\lambda_i$.   On $R_{i,j,\epsilon}$, define $\lambda_i(s,t) = \overline{F}_{i,j}(\gamma(s,t)$.   Arguing as above, it is clear that as a composition of continuous functions, $\lambda_i$ is continuous on $R_{i,j,\epsilon}$.   The first issue we have to resolve is that this definition does not conflict with the previous definiton of $\lambda_i$.   In particular, by the induction hypothesis $\lambda_i(s,t) = \overline{F}_{i,j-1}(\gamma(s,t))$ on $R_{i,j-1,\epsilon}$, so we need to make sure that the new definition agrees with the old one on $R_{i,j-1,\epsilon} \cap R_{i,j,\epsilon}$.   This follows because $\gamma(R_{i,j-1,\epsilon}\cap R_{i,j,\epsilon}) \subset U_{i,j-1} \cap U_{i,j}$ and we have defined $\overline{F}_{i,j}$ so that it equals $\overline{F}_{i,j-1}$ on $U_{i,j-1} \cap U_{i,j}$.   If $j>1$, then we observe that by Lemma \ref{banach:PrimitiveFollowingMappingThickening} $R_{i,k,\epsilon}\cap R_{i,j,\epsilon} = \emptyset$ for all $k \leq j-2$ and so there is no possible conflict from those parts of the domain of $\lambda_i$.

Having defined the $\lambda_i$ and the $\overline{F}_{s,t}$ on strips, we need to match them up to get a single definition over the entire rectangle.   Another induction argument is required.   To begin the induction, define $\lambda$ on $[s_0, s_1+\epsilon) \times [c,d]$ to be equal to $\lambda_0$.  Furthermore for $(s,t) \in [s_0, s_1] \times [c,d]$ define $F_{s,t} = \overline{F}_{s,t}$.   Then $\lambda$ is continuous, $F_{s,t}$ is holomorphic, $F^\prime_{s,t} = f$ on $U_{s,t}$ and $F_{s,t}(\gamma(u,v)) = \lambda(u,v)$ for $(u,v) \in (s - \epsilon, s + \epsilon) \times (t - \epsilon, t + \epsilon) \cap [a,b] \times [c,d]$.
  
  For the induction step, assume that $\lambda$ is defined and continuous on $[s_0, s_i+\epsilon) \times [c,d]$.
  \begin{clm}There exists a constant $C \in X$ such that $\lambda(s, t) = \lambda_i(s, t)+C$ for all $(s,t) \in (s_i-\epsilon, s_i + \epsilon) \times [c,d]$.
  \end{clm}
  By the induction hypothesis, for every $c \leq t \leq d$ we have an
  open set $U_{s_i, t}$ and a holomorphic function $F_{s_i,t} : U_{s_i,t} \to X$ such that
  \begin{itemize}
  \item[(i)] $F^\prime_{s_i,t}(z) = f(z)$ for all  $z \in U_{s_i,t}$
  \item[(ii)] $\gamma((s_i - \epsilon, s_i+ \epsilon) \times (t - \epsilon,    t+\epsilon)) \subset U_{s_i,t}$ 
  \item[(iii)]  $\lambda(u,v) = F_{s_i,t}(\gamma(u,v))$  for all $(u,v) \in    (s_i - \epsilon, s_i+ \epsilon) \times (t - \epsilon,  t+\epsilon)$
  \end{itemize}
 In addition by the construction on strip $[s_i,s_{i+1}] \times [c,d]$, there is an open set $U^i_{s_i,t}$ and a holomorphic function $\overline{F}_{s_i,t} :
 U^i_{s_i,t} \to X$ such that
 \begin{itemize}
 \item[(i)]$\gamma(s_i,t) \in U^i_{s_i,t}$
 \item[(ii)] $\lambda_i(u,v) = \overline{F}_{s_i,t}(\gamma(u,v))$ for all $(u,v) \in (s_i - \epsilon, s_i + \epsilon, t-\epsilon, t+ \epsilon)$
 \item[(iii)] $\overline{F}^\prime_{s_i,t}(z) = f(z)$ for all  $z \in U^i_{s_i,t}$
 \end{itemize}
 Since $(s_i,t) \in (s_i - \epsilon, s_i+ \epsilon) \times (t - \epsilon, t+\epsilon)$
  we know that $U_{s_i,t} \cap U^i_{s_i,t} \neq \emptyset$ and $F^\prime_{s_i,t} - \overline{F}^\prime_{s_i,t} = 0$ on $U_{s_i,t} \cap U^i_{s_i,t}$ which is connected so $F_{s_i,t} - \overline{F}_{s_i,t}$ is constant on $U_{s_i,t} \cap U^i_{s_i,t}$.   This implies that $\lambda - \lambda_i = \left(F_{s_i,t} - \overline{F}_{s_i,t}\right) \circ \gamma$ is locally constant on $(s_i-\epsilon, s_i + \epsilon) \times [c,d]$.  Since $ (s_i-\epsilon, s_i + \epsilon) \times [c,d]$ is connected and $\lambda- \lambda_i$ is continuous, this implies $\lambda - \lambda_i$ is constant on $(s_i - \epsilon, s_i + \epsilon) \times [c,d]$.
  
  Define $\lambda : [s_0, s_{i+1}+\epsilon) \times [c,d]$ by 
  \begin{align*}
    \lambda(s,t) &= \begin{cases}
      \lambda(s,t) & \text{for $s_0 \leq s \leq s_i$} \\
      \lambda_i(s,t) + C & \text{for $s_i < s \leq s_{i+1}+\epsilon$}
      \end{cases}
  \end{align*}
  By the previous claim, this definition is consistent with the existing definition of $\lambda$ on the shared domain $(s_i-\epsilon, s_i + \epsilon) \times [c,d]$
  and moreover since $\lambda$ and $\lambda_i$ are equal on $\lbrace s_i \rbrace \times [c,d]$ it follows that $\lambda$ is continuous.
  For $(s,t) \in (s_i,s_{i+1}] \times [c,d]$ define $U_{s,t} =  U^i_{s,t}$ and $F_{s,t} = \overline{F}_{s,t}+C$.   By the construction on strips, $\gamma((s - \epsilon, s + \epsilon) \times (t - \epsilon, t+\epsilon) \cap [a,b] \times [c,d]) \subset U_{s,t}$ and $F_{s,t}$ is holomorphic on $U_{s,t}$.  Since $\overline{F}_{s,t}(\gamma(u,v)) = \lambda_i(u,v)$ for every $(u,v) \in (s - \epsilon, s+\epsilon) \times (t - \epsilon, t+ \epsilon) \cap [a,b] \times [c,d]$ and $\lambda(u,v) = \lambda_i(u,v)$ for every $(u,v) \in (s_i-\epsilon, s_i + \epsilon) \times [c,d]$ we see that $F_{s,t}(\gamma(u,v)) = \lambda(u,v)$ for every $(u,v) \in (s_i-\epsilon, s_i + \epsilon) \times [c,d]$.
\end{proof}

\begin{thm}\label{banach:ContourIntegralHomotopyTheorem}Let $X$ be a complex Banach space, $U \subset \complexes$ be an open set, let $f: U \to X$ be continuous such that $\oint_\gamma f = 0$ for every rectangle $\gamma$ contained entirely in $U$ and let  $\gamma_0 : [a,b] \to U$ and $\gamma_1: [a,b] \to U$ be continuous homotopic curves, then
  \begin{align*}
   \oint_{\gamma_0} f &= \oint_{\gamma_1} f
  \end{align*}
\end{thm}
\begin{proof}
  Let $H : [a,b] \times [0,1] \to U$ be a homotopy of $\gamma_0$ and
  $\gamma_1$.   Apply Proposition
  \ref{banach:HolomorphicPrimitiveOnRectangle} to get continuous $\lambda : [a,b] \times [0,1] \to X$.  As noted in the proof of Proposition \ref{banach:HolomorphicPrimitiveOnRectangle}, $\lambda_s : [a,b] \to X$ defined by $\lambda_s(t) = \lambda(t,s)$ is a holomorphic primitive along $\gamma_s(t) = \gamma(t,s)$ and therefore by definition we have $\oint_{\gamma_0} f = \lambda(b,0) - \lambda(a, 0)$ and $\oint_{\gamma_1} f = \lambda(b,1) - \lambda(a, 1)$.   On the other hand, we also know that $\lambda_a(s) = \lambda(a,s)$ (respectively $\lambda_b(s) = \lambda(b,s)$) is a holomorphic primitive of $f$ along $\gamma_a(s) = H(a,s) = \gamma_0(a) = \gamma_1(a)$ (respectively along $\gamma_b(s) = H(b,s) = \gamma_0(b) = \gamma_1(b)$).   Therefore by Corollary \ref{banach:PointContourIntegralZero}  and definition
  \begin{align*}
    0 &= \oint_{\gamma_a} = \lambda_a(1) - \lambda_a(0) = \lambda(a, 1) - \lambda(a, 0) \\
    0 &= \oint_{\gamma_b} = \lambda_b(1) - \lambda_b(0) = \lambda(b, 1) - \lambda(b, 0) \\
  \end{align*}
  which shows
  \begin{align*}
    \oint_{\gamma_0} f &= \lambda(b,0) - \lambda(a, 0) = \lambda(b,1) - \lambda(a, 1)  = \oint_{\gamma_1} f
  \end{align*}
\end{proof}

\begin{thm}[Cauchy's Theorem]\label{banach:CauchysTheorem}Let $X$ be a complex Banach space, $U \subset \complexes$ be an open set, let $f: U \to X$ be continuous such that $\oint_\gamma f = 0$ for every rectangle $\gamma$ contained entirely in $U$ and let  $\gamma : [a,b] \to U$ and let $\gamma$ be homotopic to a point in $U$ then
  \begin{align*}
   \oint_{\gamma} f &= 0
  \end{align*}
\end{thm}
\begin{proof}
  By Corollary \ref{banach:PointContourIntegralZero} we have $\oint_{\tilde{\gamma}} f = 0$ for any point contour $\tilde{\gamma}$, so this result follows immediately from Theorem \ref{banach:ContourIntegralHomotopyTheorem}.
\end{proof}

\begin{defn}Let $U \subset \complexes$ be a open set, $\gamma : [a,b] \to U$ and $z_0 \in U$ such that $z_0$ is not contained in the image of $\gamma$ the we let
  \begin{align*}
    I(\gamma, z_0) &= \frac{1}{2\pi i}\oint_\gamma \frac{dz}{z - z_0}
  \end{align*}
  denote the \emph{index of $\gamma$ with respect to $z_0$}.
\end{defn}

\begin{thm}[Cauchy's Integral Formula]\label{banach:CauchysIntegralFormula}Let $X$ be a complex Banach space, $U \subset \complexes$ be an open set, let $f: U \to X$ be holomorphic and let  $\gamma : [a,b] \to U$ and let $\gamma$ be homotopic to a point in $U$ then for every $z_0 \in U$ such that $z_0$ is not contained in the image of $\gamma$,
  \begin{align*}
   f(z_0) \cdot I(\gamma, z_0) &= \frac{1}{2\pi i} \oint_\gamma \frac{f(z)}{z - z_0} \, dz
  \end{align*}
\end{thm}
\begin{proof}
  Define the function
  \begin{align*}
    g(z) &= \begin{cases}
      \frac{f(z) - f(z_0)}{z-z_0} & \text{ for $z \in U \setminus \lbrace z_0 \rbrace$} \\
      f^\prime(z_0) & \text{ for $z = z_0$}
    \end{cases}
  \end{align*}
  Since $f$ is holomorphic, it follows that $g$ is holomorphic on $U \setminus \lbrace z_0 \rbrace$ (hence continuous on $U \setminus \lbrace z_0 \rbrace$).  Also since $f$ is holomorphic, we know that $\lim_{z \to z_0} g(z) = f^\prime(z_0) = g(z_0)$, so $g$ is continuous on all of $U$.   By the punctured Goursat Lemma \ref{banach:GoursatLemmaPunctured} we conclude that $\oint_R g = 0$ for all rectangles entirely contained in $U$ which in turn allows us to conclude by Cauchy's Theorem \ref{banach:CauchysTheorem} that $\oint_\gamma g = 0$.  Applying the definition of $g$ and the index,
  \begin{align*}
    0 &= \oint_\gamma g = \oint_\gamma  \frac{f(z) - f(z_0)}{z-z_0}  \, dz = \oint_\gamma  \frac{f(z)}{z-z_0}  \, dz - \oint_\gamma  \frac{f(z_0)}{z-z_0}  \, dz = \oint_\gamma  \frac{f(z)}{z-z_0}  \, dz - 2 \pi i I(\gamma, z_0)
  \end{align*}
  and the result follows.
\end{proof}

TODO: Show index of a circle equals 1, show that index is always an integer, show every holomorphic function is analytic

We are just about ready to show that every holomorphic function is complex analytic, but we need quick lemma about the convergence of the Neumann series in the complex domain.

\begin{lem}\label{banach:ComplexNeumann}The power series $\sum_{n=0}^\infty z^n$ converges absolutely and uniformly on any disk $\abs{z} \leq \rho < 1$.   For any such $\abs{z} < 1$,
  \begin{align*}
    \sum_{n=0}^\infty z^n &= \frac{1}{1 - z}
  \end{align*}
\end{lem}
\begin{proof}
  With $\rho < 1$ fixed and working in $\sboundedfunctions{\overline{B(0, \rho)}, X}$, we note that the function $\norm{z^n}_\infty \leq \rho^n$ for each $n \geq 0$.   Since $\rho < 1$, we have $\sum_{n=0}^\infty \rho^n < \infty$ and therefore the result follows by the Weierstrass M Test (Theorem \ref{banach:WeierstrassMTest}).
\end{proof}

\begin{thm}\label{banach:ComplexAnalyticHolomorphicEquivalence}Let $X$ be a complex Banach space, $U \subset \complexes$ be an open set and let $f : U \to X$ then $f$ is complex analytic on $U$ if and only if $f$ is holomorphic on $U$.  More specifically, if $f$ is holomorphic on a disk $B(z_0, r_0)$ then it has a convergent power series on $B(z_0, r_0)$.
\end{thm}
\begin{proof}
  We have already seen in Proposition \ref{banach:ComplexAnalyticImpliesHolomorphic} that a complex analytic function $f$ on $U$ is holomorphic on $U$.  Suppose that $f$ is holomorphic on a disk $B(z_0, r)$.   Let $0 < r < r_0$ be given and we will show that there is a power series which converges to $f$ absolutely and uniformly on $\abs{z - z_0} \leq r$.   Pick $0 < r < \rho < r_0$ and let $\abs{z - z_0} = \rho$ denote the standard circle of radius $\rho$ centered at $z_0$.   Then by the Cauchy Integral Formula \ref{banach:CauchysIntegralFormula} we have
  \begin{align*}
   f(w) &= \frac{1}{2\pi i}  \oint_{\abs{z - z_0} = \rho} \frac{f(z)}{z - w} \, dz \text{ for all $\abs{w - z_0} \leq r$}
  \end{align*}

  Next we develop $\frac{1}{z - w}$ in a power series based on the Neumann series where we use the fact that $\abs{\frac{w - z_0}{z - z_0}} \leq \frac{r}{\rho} < 1$ so by Lemma \ref{banach:ComplexNeumann}:
  \begin{align*}
    \frac{1}{z - w} &=   \frac{1}{z - z_0 - w + z_0}  = \frac{1}{z - z_0} \cdot \frac{1}{1 - \frac{w - z_0}{z - z_0}} \\
    &=\frac{1}{z - z_0} \sum_{n=0}^\infty \left(\frac{w - z_0}{z - z_0}\right)^n
  \end{align*}
  Also by Lemma \ref{banach:ComplexNeumann}, the series $\sum_{n=0}^\infty \left(\frac{w - z_0}{z - z_0}\right)^n$ converges absolutely and uniformly over $\abs{z - z_0} = \rho$ (with $w$ fixed).   Moreover, $\frac{f(z)}{z - z_0}$ is holomorphic hence continuous on $\abs{z - z_0} = \rho$ and therefore bounded by compactness (Heine-Borel Theorem).

  \begin{clm}For fixed $\abs{w} \leq r$, $\sum_{n=0}^\infty \frac{f(z)}{z - z_0}\left(\frac{w - z_0}{z - z_0}\right)^n$ converges absolutely and uniformly on $\abs{z - z_0} = \rho$.
  \end{clm}
  TODO: Proof of claim


  Now with the claim in hand we note that we can exchange the order of summation and integration,
  for by Proposition \ref{banach:ArclengthProperties},
  \begin{align*}
    \norm{\oint_\gamma f_n - \oint_\gamma f} &= \norm{\oint_\gamma (f_n -f) } \leq \sup_{z \in \gamma} \norm{f(z) - f_n(z)} \abs{\gamma}
  \end{align*}
  so clearly if $\lim_{n \to \infty} \sup_{z \in \gamma} \norm{f_n(z) - f(z)} = 0$ then $\lim_{n \to \infty} \oint_\gamma f_n = \oint_\gamma f$.   This extends easily to series by treating the sequence of partial sums and using linearity of the integral
  \begin{align*}
    \oint_\gamma \sum_{n=0}^\infty f_n &= \oint_\gamma \lim_{m \to \infty} \sum_{n=0}^m f_n = \lim_{m \to \infty} \oint_\gamma \sum_{n=0}^m f_n \\
    &=\lim_{m \to \infty} \sum_{n=0}^m \oint_\gamma f_n = \sum_{n=0}^\infty \oint_\gamma f_n 
  \end{align*}
So applied in this case
  \begin{align*}
    f(w) &= \frac{1}{2\pi i} \oint_{\abs{z - z_0} = \rho} \frac{f(z)}{z - w} \, dz \\
    &=\frac{1}{2\pi i}  \oint_{\abs{z - z_0} = \rho} \sum_{n=0}^\infty \frac{f(z)}{z - z_0}\left(\frac{w - z_0}{z - z_0}\right)^n \, dz \\
    &= \sum_{n=0}^\infty \frac{(w - z_0)^n}{2\pi i} \oint_{\abs{z - z_0} = \rho} \frac{f(z)}{(z - z_0)^{n+1}} \, dz \\
  \end{align*}
\end{proof}

\begin{thm}[Liouville's Theorem]\label{banach:LiouvillesTheorem}Let $X$ be a complex Banach space and $f : \complexes \to X$ be a bounded entire function then $f$ is constant.
\end{thm}
\begin{proof}
Since $f$ is entire we may write $f(z) = \sum_{n=0}^\infty v_n z^n$ where by Cauchy's Integral Formula \ref{banach:CauchysIntegralFormula} we have for every $0 < r < \infty$,
\begin{align*}
v_n &= \frac{1}{2\pi i} \oint_{\abs{z} = r} \frac{f(z)}{z^{n+1}} \, dz
\end{align*}
Suppose that $f$ is bounded so that $\norm{f(z)} \leq M$ for all $z \in \complexes$ then
\begin{align*}
\norm{v_n} &\leq \frac{1}{2\pi} \oint_{\abs{z} = r} \frac{\norm{f(z)}}{\abs{z}^{n+1}} \, d\abs{z} \leq \frac{1}{2\pi} \oint_{\abs{z} = r} \frac{M}{r^{n+1}} \, d\abs{z} \frac{M}{2\pi r^{n+1}}  2 \pi r = \frac{M}{r^n}
\end{align*}
Since $0 < r < \infty$ was arbitrary we let $r \to \infty$ and observe that $v_n = 0$ for every $n > 0$.
\end{proof}

